% Part: turing-machines
% Chapter: machines-computations
% Section: combining-machines

\documentclass[../../../include/open-logic-section]{subfiles}

\begin{document}

\olfileid{tur}{mac}{cmb}
\olsection{ટ્યુરિંગ મશીનોનું સંયોજન}

\begin{explain}
અત્યાર સુધી આપણે જોયેલાં ટ્યુરિંગ મશીનો સ્વરૂપે ઘણાં
સરળ હતાં. પરંતુ હકીકતમાં, કોઈ પણ આધુનિક પ્રોગ્રામિંગ
ભાષામાં ઉકેલી શકાય તેવી સમસ્યા ટ્યુરિંગ મશીનોથી પણ
ઉકેલી શકાય છે. વધુ જટિલ ટ્યુરિંગ મશીનો બનાવવા માટે
તેમને જોડવાનું શક્ય છે તે સમજવું જરૂરી છે. એમ હોય તો
પ્રક્રિયાને સરળ ભાગોમાં વહેંચીને વધુ જટિલ સમસ્યાઓ
ઉકેલતાં મશીનો બનાવી શકીએ. જટિલ સમસ્યાને તેના
ઘટક ભાગોમાં વહેંચવાની સ્વાભાવિક રીત મળે, તો અનેક
સરળ ટ્યુરિંગ મશીનો બનાવીને તેમને એક મશીનમાં જોડતાં
સમસ્યાને કેટલાંક તબક્કામાં ઉકેલી શકીએ. આગળના વિભાગમાં
અટકવાની સમસ્યા હાથ ધરતી વખતે આ વાત ખાસ અગત્યની છે.

ટ્યુરિંગ મશીનો $M = \tuple{Q, \Sigma, q_0, \delta}$
અને~$M' = \tuple{Q', \Sigma', q_0', \delta'}$ને કેવી રીતે
જોડીએ? હવે $M$ થંભે તે પછી ટેપનું સંરૂપણ મશીન~$M'$
ના ચલનના ઇનપુટ સંરૂપણ તરીકે વાપરીએ. આ કામ કરતું
એક ટ્યુરિંગ મશીન $M \frown M'$ મેળવવા માટે આવું કરો:
\begin{enumerate}
    \item $M'$ની બધી અવસ્થાઓ~$Q'$ના ક્રમાંક (અથવા નામ)
    ફરી ગોઠવો, જેથી $M$ અને~$M'$ની કોઈ અવસ્થા
    સમાન ન હોય ($Q \cap Q' = \emptyset$).
    \item $M \frown M'$ની અવસ્થાઓ $Q \cup Q'$ છે.
    \item ટેપની વર્ણમાળા $\Sigma \cup \Sigma'$ છે.
    \item પ્રારંભિક અવસ્થા~$q_0$ છે.
    \item સંક્રમણ વિધેય નીચે આપેલું વિધેય $\delta''$ છે:
    \[\delta''(q,\sigma) =
    \begin{cases}
      \delta(q,\sigma) & \text{જો $q \in Q$ અને $\delta(q,\sigma)$ વ્યાખ્યાયિત હોય}\\
      \delta'(q,\sigma) & \text{જો $q \in Q'$ હોય}\\
      \tuple{q_0', \sigma, \TMstay} & \text{જો $q \in Q$ હોય અને
      $\delta(q,\sigma)$ અવ્યાખ્યાયિત હોય}
    \end{cases}\]
\end{enumerate}
મળેલું મશીન અવસ્થા $q \in Q$માં હોય ત્યારે~$M$ની
સૂચનાઓ અને અવસ્થા~$q \in Q'$માં હોય ત્યારે~$M'$
ની સૂચનાઓ વાપરે છે. તે અવસ્થા $q \in Q$માં હોય
અને એવું પ્રતીક~$\sigma$ વાંચે કે જેના માટે $M$માં
સંક્રમણ નથી (એટલે $M$ થંભ્યું હોત), ત્યારે તે~$M'$
ની પ્રારંભિક અવસ્થામાં જાય છે (અને ટેપની સામગ્રી તથા
હેડનું સ્થાન યથાવત્ રાખે છે).
\sourcecorrection{OLTM-004}{સ્થિર અંગ્રેજી મૂળના પ્રથમ
કિસ્સામાં માત્ર $q\in Q$ લખ્યું છે, તેથી $\delta(q,\sigma)$
અવ્યાખ્યાયિત હોય ત્યારે પણ પહેલો કિસ્સો લાગુ પડે અને
ત્રીજો કિસ્સો ઢંકાઈ જાય. ગુજરાતી સૂત્રના પહેલા કિસ્સામાં
$\delta(q,\sigma)$ વ્યાખ્યાયિત હોવાની શરત સ્પષ્ટ કરી છે;
થંભેલા $M$માંથી $M'$માં જવાનું બાજુનું વર્ણન આ અર્થ
નક્કી કરે છે.}

નોંધો કે મશીન~$M$ શિસ્તબદ્ધ ન હોય, તો $M$ થંભે
ત્યારે ટેપનો હેડ ક્યાં હશે તે આપણે જાણતા નથી; તેથી
$M$ના થંભવાના સંરૂપણમાં હેડ ખાનું~$1$ વાંચતો હોય
એવું જરૂરી નથી. મશીનો જોડતી વખતે આ વાત ધ્યાનમાં
રાખવી અગત્યની છે.
\end{explain}

\begin{ex}
\emph{મશીનોનું સંયોજન:} એવું મશીન રચીશું કે તેને ઇનપુટ
તરીકે $n$ અને~$m$ લંબાઈના $\TMstroke$ના બે સમૂહો
આપીએ, તો તે ટેપ પર $2(m+n)$ $\TMstroke$ના એક
જ સમૂહ સાથે થંભે. તેને બનાવવા આપણે પહેલેથી જાણીએ
છીએ એવાં બે મશીનો---સરવાળાનું મશીન અને બમણું કરનાર
મશીન---જોડી શકીએ. પહેલાં સરવાળાના મશીનનો અવસ્થા-આલેખ
દોરીએ.
\[
\begin{tikzpicture}[->,>=stealth',shorten >=1pt,auto,node distance=2.8cm,
                    semithick]
  \tikzstyle{every state}=[fill=none,draw=black,text=black]

  \node[initial,state] (A)              {$q_0$};
  \node[state]         (B) [right of=A] {$q_1$};
  \node[state]         (C) [right of=B] {$q_2$};

  \path (A) edge node {\TMtrans{\TMblank}{\TMstroke}{\TMstay}} (B)
            edge [loop above] node {\TMtrans{\TMstroke}{\TMstroke}{\TMright}} (B)
        (B) edge [loop above] node {\TMtrans{\TMstroke}{\TMstroke}{\TMright}} (B)
            edge node {\TMtrans{\TMblank}{\TMblank}{\TMleft}} (C)
        (C) edge [loop above] node {\TMtrans{\TMstroke}{\TMblank}{\TMstay}} (C);
\end{tikzpicture}
\]
અવસ્થા~$q_2$માં થંભવાને બદલે, આઉટપુટ બમણું કરવા
મશીન આગળ ચાલે એવું જોઈએ. યાદ કરો કે બમણું કરનાર
મશીન ઇનપુટનો પ્રથમ સ્ટ્રોક ભૂંસીને અલગ આઉટપુટમાં
બે સ્ટ્રોક લખે છે. ટેપનો હેડ સરવાળાના મશીનના
આઉટપુટનો પ્રથમ સ્ટ્રોક વાંચે તેની ખાતરી કરતી સૂચના
ઉમેરીએ.
\[
\begin{tikzpicture}[->,>=stealth',shorten >=1pt,auto,node distance=2.8cm,
                    semithick]
  \tikzstyle{every state}=[fill=none,draw=black,text=black]

  \node[initial,state] (A)              {$q_0$};
  \node[state]         (B) [right of=A] {$q_1$};
  \node[state]         (C) [right of=B] {$q_2$};
  \node[state]         (D) [below left of=C] {$q_3$};
  \node[state]         (E) [below left of=D] {$q_4$};

  \path (A) edge node {\TMtrans{\TMblank}{\TMstroke}{\TMstay}} (B)
            edge [loop above] node {\TMtrans{\TMstroke}{\TMstroke}{\TMright}} (B)
        (B) edge [loop above] node {\TMtrans{\TMstroke}{\TMstroke}{\TMright}} (B)
            edge node {\TMtrans{\TMblank}{\TMblank}{\TMleft}} (C)
        (C) edge node {\TMtrans{\TMstroke}{\TMblank}{\TMleft}} (D)
        (D) edge [loop left] node {\TMtrans{\TMstroke}{\TMstroke}{\TMleft}} (D)
            edge node[left, xshift=-2mm] {\TMtrans{\TMendtape}{\TMendtape}{\TMright}} (E);
\end{tikzpicture}
\]
હવે ઇનપુટ બમણું કરવું સહેલું છે---બમણું કરનાર મશીનને
અવસ્થા~$q_4$ સાથે જોડવાનું જ બાકી છે. એ માટે
બમણું કરનાર મશીનની અવસ્થાઓનાં નામ બદલવાં પડે,
જેથી તે~$q_0$ને બદલે~$q_4$થી શરૂ થાય---આ રીતે
બે પ્રારંભિક અવસ્થાઓ ઊભી થતી નથી. આખરી આલેખ
\olref{fig:combined} જેવો હોવો જોઈએ.
\begin{figure}
\[
\begin{tikzpicture}[->,>=stealth',shorten >=1pt,auto,node distance=2.8cm,
                    semithick]
  \tikzstyle{every state}=[fill=none,draw=black,text=black]
  \node[initial,state] (A)              {$q_0$};
  \node[state]         (B) [right of=A] {$q_1$};
  \node[state]         (C) [right of=B] {$q_2$};
  \node[state]         (D) [below left of=C] {$q_3$};
  \node[state]         (E) [below left of=D] {$q_4$};
  \node[state]         (2) [right of=E] {$q_5$};
  \node[state]         (3) [right of=2] {$q_6$};
  \node[state]         (4) [below of=3] {$q_7$};
  \node[state]         (5) [left of=4]  {$q_8$};
  \node[state]         (6) [left of=5]  {$q_9$};

  \path (A) edge node {\TMtrans{\TMblank}{\TMstroke}{\TMstay}} (B)
            edge [loop above] node {\TMtrans{\TMstroke}{\TMstroke}{\TMright}} (B)
        (B) edge [loop above] node {\TMtrans{\TMstroke}{\TMstroke}{\TMright}} (B)
            edge node {\TMtrans{\TMblank}{\TMblank}{\TMleft}} (C)
        (C) edge node {\TMtrans{\TMstroke}{\TMblank}{\TMleft}} (D)
        (D) edge [loop left] node {\TMtrans{\TMstroke}{\TMstroke}{\TMleft}} (D)
            edge node[left, xshift=-2mm] {\TMtrans{\TMendtape}{\TMendtape}{\TMright}} (E)
    (E) edge node {\TMtrans{\TMstroke}{\TMblank}{\TMright}} (2)
    (2) edge [loop above] node {\TMtrans{\TMstroke}{\TMstroke}{\TMright}} (2)
      edge node {\TMtrans{\TMblank}{\TMblank}{\TMright}} (3)
    (3) edge [loop above] node {\TMtrans{\TMstroke}{\TMstroke}{\TMright}} (3)
        edge node {\TMtrans{\TMblank}{\TMstroke}{\TMright}} (4)
    (4) edge [loop below] node {\TMtrans{\TMblank}{\TMstroke}{\TMleft}} (4)
        edge node {\TMtrans{\TMstroke}{\TMstroke}{\TMleft}} (5)
    (5) edge [loop below]  node {\TMtrans{\TMstroke}{\TMstroke}{\TMleft}} (5)
        edge              node {\TMtrans{\TMblank}{\TMblank}{\TMleft}} (6)
    (6) edge [loop below] node {\TMtrans{\TMstroke}{\TMstroke}{\TMleft}} (6)
        edge              node {\TMtrans{\TMblank}{\TMblank}{\TMright}} (E);
\end{tikzpicture}
\]
\caption{સરવાળાનું અને બમણું કરનાર મશીન જોડવાં}
\ollabel{fig:combined}
\end{figure}
\end{ex}

\begin{prop}
    $M$ અને $M'$ શિસ્તબદ્ધ હોય અને અનુક્રમે
    $f\colon \Nat^k \to \Nat$ અને
    $f'\colon \Nat \to \Nat$ વિધેયો ગણતા હોય, તો
    $M \frown M'$ પણ શિસ્તબદ્ધ છે અને
    $\comp{f}{f'}$ ગણે છે.
\end{prop}

\begin{proof}
    $M$ શિસ્તબદ્ધ હોવાથી તે
    આઉટપુટ~$f(n_1,\dots,n_k) = m$ સાથે થંભે ત્યારે
    હેડ ખાનું~$1$ વાંચે છે. હવે~$M'$ની પ્રારંભિક
    અવસ્થામાં જઈએ, તો મશીન આઉટપુટ $f'(m)$ સાથે,
    ફરી ખાનું~$1$ વાંચતાં, થંભશે.
    \olref[dis]{defn:disciplined}ની બીજી શરતો પણ
    પૂરી થાય છે.
\end{proof}

\begin{prob}
    \cref{tur:mac:dis:prob:disc-succ}નું મશીન~$M$
    લઈને અને $M \frown M$ રચીને $f(x) = x+2$
    ગણતું શિસ્તબદ્ધ ટ્યુરિંગ મશીન આપો.
\end{prob}
\end{document}
